\section{Grounding under a Token Budget: a Diagnosis}
\label{sec:diagnosis}

We first establish Observations 1 and 2 of the introduction on S-EMBER, before introducing any method.
Observation 3, which concerns the answer format, is studied in \cref{sec:transfer}.
Unless stated otherwise, the reference memory in this section is random 10\% pruning.

\subsection{Setting}
\label{sec:setting}

A question $x$ is asked at time $q$ about a video $v_{0:q}$. The model sees $F{=}64$ frames from $[0,q]$, which are
encoded into a KV cache $\Mem$ of $N$ visual tokens. Timestamp tokens are interleaved with the visual tokens, and
Qwen3-VL uses one ``\texttt{<t seconds>}'' for each pair of frames. A budget $\rho$ (by default $\rho{=}0.1$) limits
the number of tokens kept in each layer to $\rho N$, and kept timestamps count toward this budget. The model answers
with a short \emph{timeline}:
\begin{center}
\small\ttfamily
\begin{adjustbox}{max width=\linewidth}\begin{tabular}{l}
Seen: $t_1$ seconds, \dots \\
Seen: $t_k$ seconds, \dots \\
Answer: \dots \\
Time: [$s$, $e$]
\end{tabular}\end{adjustbox}
\end{center}
and we parse the interval $[s,e]$ from the last line. This reasoning-style format suits grounded QA, where the answer
describes what was seen before it gives the interval, and we use it throughout the diagnosis. \Cref{sec:transfer}
compares it with the one-line format of grounding benchmarks (\emph{``The event happens in $s$--$e$ seconds''}).
Grounding is scored by IoU with the gold interval, reported as mIoU and R@$\{0.3, 0.5, 0.7\}$. Grounded QA also
requires the answer to be judged correct, so GQ@0.5 counts answers that are correct and have IoU $\geq 0.5$. We call
an answer an \emph{early start} when its interval starts in the first 5\% of the window, and a \emph{leak} when it
starts early even though the gold evidence starts later.

\subsection{QA pruners lose temporal coverage}
\label{sec:coverage-loss}

At 10\% of the tokens, FastVID, VidCom\textsuperscript{2}, FlashVID and SemVID all score below random token
selection (\cref{sec:pruners}). Where the kept tokens lie explains why (\cref{tab:retention}). HERMES scores tokens by
attention plus recency, so it keeps 27.6\% of the tokens in the last quarter of the window and only 2.5--4.7\%
elsewhere. FastVID retains one representative per run of similar frames, and as a result it keeps a token from only
25\% of the sampled frames in a median question. Those runs of similar frames are exactly where a continuing event
lives. Random selection keeps about 9\% everywhere and never empties a frame. Because a timestamp is useful only when
some visual content of its frame survives with it, a selector that empties frames leaves holes in the timeline. These
holes hurt even when the score itself is sensible. If we rank HERMES's scores within each frame instead of globally,
which gives every frame an equal share (``stratified''), mIoU improves by $5.3$.

\begin{table}[t]
  \centering
  \small
  \setlength{\tabcolsep}{3.5pt}
  \caption{\textbf{Where pruning keeps tokens.} Share of the visual tokens kept, by position in the window
  (S-EMBER, 10\% budget), and the resulting grounding (mIoU) and leak. FastVID (mIoU 17.9, leak 38\%) keeps
  6.0\,/\,8.5\,/\,10.5\,/\,10.6\,/\,12.9\% over the five fifths of the window, and tokens from only 25\% of the frames.}
  \label{tab:retention}
  \begin{adjustbox}{max width=\linewidth}\begin{tabular}{@{}lccccccc@{}}
    \toprule
    Selector & 0--10\% & 10--25\% & 25--50\% & 50--75\% & 75--100\% & mIoU & Leak \\
    \midrule
    HERMES     & 4.7 & 3.9 & 2.6 & 2.5 & 27.6 & 18.3 & 49\% \\
    Stratified & 8.7 & 8.7 & 8.6 & 8.6 & 8.7 & \second{23.7} & \best{29\%} \\
    Random     & 8.9 & 8.8 & 8.6 & 8.7 & 8.6 & \best{25.8} & \second{30\%} \\
    \bottomrule
  \end{tabular}\end{adjustbox}
\end{table}

\subsection{The prior fills the gaps}
\label{sec:prior}

\paragraph{The model answers without video.}
When we ask every question with an empty visual memory, the model still writes a timeline. Of its first
\texttt{Seen:} moments, 89\% are at 0\,s, and 90\% of its intervals start in the first 5\% of the window, which gives
an 88\% leak. The median interval is 0.45$\times$ the gold length, and 87\% of the ends fall more than half an
interval too early. We call this habit of starting early and stopping soon the model's \emph{temporal prior}.

\paragraph{The prior leaks into answers with video.}
The same signature appears, more weakly, in every memory (\cref{fig:leak}). About 30\% of intervals start early at
every random budget and also with no pruning at all, compared with 16\% of gold intervals. The share rises when
coverage is lost, to 49\% for HERMES and 37--39\% for HERMES in streaming mode. The leak therefore depends on where the
memory is silent and not on how many tokens it keeps.

\paragraph{Removing the evidence causes the leak.}
To test causality, we drop the uniform frames that fall inside the gold interval (``evidence''), or the same number
of frames elsewhere (``control''), and keep every other frame with its true timestamp (\cref{tab:causal}). Without
the evidence, the answers do not spread out. They move to the start of the video instead, and the leak rises by 14
points for both the unpruned and the random 10\% memory. In the HERMES 10\% memory, the leak also tracks how much
of the evidence survives. Questions whose evidence keeps under 5\% of its tokens leak 56\% of the time, compared with 33\% when at least 15\% survives.

\paragraph{An instruction does not remove it.}
We added ``The relevant moments may occur at any point in the video; do not assume they are near the start'' to the
prompt. The leak barely moved (29.5\%~$\rightarrow$~28.0\%), and grounding did not change ($-0.1$ mIoU). In this
answer format, then, the prior acts in the output distribution and has to be removed there. We do not yet know
whether the same prior appears when the model answers with a single interval. In that format PCD has no effect
(\cref{sec:transfer}). Those runs omit the answer reserve, which on its own changes S-EMBER grounding by only
$+0.1$. The null effect suggests that the prior decides fewer of the written times there. \TODO{no-video answers and leak
under the one-line and direct-answer prompts.}

\paragraph{The end is the weak boundary.}
Without contrastive decoding, the median end error across memories is 0.40--0.69 of the gold length, while the
median start error is 0.14--0.30. The most frequent failure is an interval that lies entirely \emph{before} the gold
one, which accounts for 33\% of all answers (\cref{tab:errors}).

% Early starts (interval starts in the first 5% of the window) per memory, greedy vs Contrast. S-EMBER.
\begin{figure}[t]
\centering
\begin{tikzpicture}
\begin{axis}[
    width=\linewidth, height=5.3cm, ybar, bar width=5.2pt,
    ymin=0, ymax=100, ylabel={early starts (\%)}, ylabel style={font=\sffamily\scriptsize},
    symbolic x coords={novid,all,r5,r10,r25,strat,herm,s4,s6}, xtick=data,
    xticklabels={no video, all tokens, random 5\%, random 10\%, random 25\%, stratified, HERMES 10\%,
                 stream 4k, stream 6k},
    xticklabel style={font=\sffamily\scriptsize, rotate=35, anchor=north east},
    yticklabel style={font=\sffamily\scriptsize},
    legend style={font=\sffamily\scriptsize, at={(0.98,0.97)}, anchor=north east, draw=none, fill=white,
                  fill opacity=0.85, text opacity=1},
    legend cell align=left, enlarge x limits=0.07, axis x line*=bottom, axis y line*=left,
    nodes near coords, nodes near coords style={font=\sffamily\tiny}]
  \addplot[fill=black!35, draw=black!55] coordinates
    {(novid,90) (all,41) (r5,39) (r10,38) (r25,37) (strat,36) (herm,56) (s4,45) (s6,46)};
  \addplot[fill=ours!80, draw=ours] coordinates
    {(all,23) (r5,20) (r10,19) (r25,21) (strat,17) (herm,31) (s4,24) (s6,24)};
  \draw[gain, dashed, thick] (rel axis cs:0,0.16) -- (rel axis cs:1,0.16)
    node[pos=0.995, font=\sffamily\tiny, anchor=south east, text=gain, fill=white, inner sep=1pt] {gold: 16\%};
  \legend{greedy, Contrast}
\end{axis}
\end{tikzpicture}
\caption{\textbf{The temporal prior leaks into every memory, and Contrast removes most of it.} Share of predicted
intervals that start in the first 5\% of the window (S-EMBER). Without video the model starts early 90\% of the
time; every memory inherits part of this habit, most where pruning loses coverage (HERMES). Contrast roughly halves
it everywhere.}
\label{fig:leak}
\end{figure}


\begin{table}[t]
  \centering
  \small
  \setlength{\tabcolsep}{3pt}
  \caption{\textbf{Causal test.} Frames inside the gold evidence (``evidence'') or as many frames elsewhere
  (``control'') are removed. Acc.\,/\,mIoU\,/\,GQ@0.5 and leak; leak measured on the 397 questions whose evidence starts
  after the first 5\% of the window. PCD: \cref{sec:contrast}.}
  \label{tab:causal}
  \begin{adjustbox}{max width=\linewidth}\begin{tabular}{@{}llcccc@{}}
    \toprule
    & & \multicolumn{2}{c}{Greedy} & \multicolumn{2}{c}{PCD} \\
    \cmidrule(lr){3-4} \cmidrule(l){5-6}
    Memory & Removed & Acc./mIoU/GQ & Leak & Acc./mIoU/GQ & Leak \\
    \midrule
    Random 10\% & control  & \second{16.6}/\second{28.8}/\second{8.8} & \best{31\%} & \second{15.4}/\second{34.6}/\second{10.5} & \best{11\%} \\
                & evidence & \phantom{0}7.8/\phantom{0}8.6/1.7 & 46\% & \phantom{0}9.3/10.3/1.1 & 23\% \\
    All tokens  & control  & \best{19.6}/\best{29.4}/\best{10.1} & \second{37\%} & \best{18.1}/\best{35.9}/\best{12.6} & \second{15\%} \\
                & evidence & \phantom{0}9.1/\phantom{0}9.4/1.9 & 51\% & \phantom{0}9.1/10.5/1.5 & 22\% \\
    \bottomrule
  \end{tabular}\end{adjustbox}
\end{table}
